\section{The split-prime step in the definite Siegel--Weil induction}

% BEGIN Z75D1
\paragraph{Z75D1. Degree zero for the actual class-sum identity.}
Let $E$ be a CM field, $\Psi$ a CM type, $n$ the number of rows,
and $d$ the degree (number of columns). For an admitted lattice
family $\phi$ at levels $M,R$, denote by $D_\phi(Z)$ the native
class sum of $\Theta_H(\phi_\Lambda;Z,0,0)-
\mathrm{Eis}(\phi_\Lambda;Z,0,0)$. Here $Z=(Z_\tau)_{\tau\in\Psi}$
has positive definite Hermitian imaginary parts. The assertion
$\mathrm{SWAt}(n,d)$ asks that $D_\phi(Z)=0$ for every such family
and every such period of degree $d$.

When $d=0$, the only $n$-by-$0$ argument of a test $f$ is zero.
The theta sum is $f(0)$, since its exponent is zero. The split
unitary group and its set of bottom spaces each have one element.
The empty determinant is one and the elliptic exponent vanishes.
The test $f$ itself is Schwartz--Bruhat at level one and satisfies
the defining transformation equation for the Weil transform. Every
test satisfying that equation has value $f(0)$, by evaluating its
theta sum. Thus the source's chosen transform contributes exactly
$f(0)$ in the single Eisenstein summand. The difference is zero
separately for every lattice before taking the class sum. This proves
the base case with no restriction on the values of $\phi_\Lambda(0)$.
% END Z75D1

% BEGIN Z75D2
\paragraph{Z75D2. A degree-one prime ideal suitable for every level.}
Write $O=\mathcal O_E$ and $\bar{\phantom{a}}$ for CM conjugation.
For every positive integer $m$, there exist a rational prime $q$
and a maximal ideal $P$ of $O$ satisfying
\[
 (q,m)=1,\qquad |O/P|=q,\qquad P+\bar P=O,\qquad
 q\nmid 2\operatorname{disc}(E).
\]
These are all the requirements of the source's $\mathrm{GoodPrime}(m,q,P)$.

Choose an integral primitive element $\alpha$ of $E/\mathbb Q$ by
clearing the denominator of a primitive element. The order
$\mathbb Z[\alpha]$ has finite index $I$ in $O$ because both are
full integer lattices in $E$. Since conjugation is nontrivial and
$\alpha$ generates $E$, the integer $N_{E/\mathbb Q}(\alpha-\bar\alpha)$
is nonzero. Form the positive integer
\[
 B=\left|2mI\operatorname{disc}(E)
              N_{E/\mathbb Q}(\alpha-\bar\alpha)\right|.
\]
Let $f$ be the monic minimal polynomial of $\alpha$. Choose an integer
$a$ such that $c=f(a)\ne0$. The nonconstant integer polynomial
$f(a+cBt)/c$ has the form $1+Bt\,g(t)$ with $g\in\mathbb Z[t]$.
At an integer $t$ where its absolute value exceeds one, take a prime
divisor $q$ of its value. Then $q\nmid B$ and $f$ has the root
$a+cBt$ modulo $q$.

Evaluation at that root is a surjection $\mathbb Z[\alpha]\to\mathbb F_q$.
It extends to $O$ since $IO\subset\mathbb Z[\alpha]$ and $I$ is a
unit modulo $q$ (equivalently the two orders agree after inverting $I$).
Its kernel $P$ has residue field $\mathbb F_q$. If $P=\bar P$,
conjugation induces an automorphism of this prime field, hence is
the identity on it. Then $\alpha-\bar\alpha\in P$. Multiplication
by this element modulo $q$ on $O/qO$ has image in the proper
subspace $P/qO$, so its determinant, the field norm, is zero
modulo $q$. This contradicts $q\nmid B$. The two distinct maximal
ideals are comaximal, as required.
% END Z75D2

% BEGIN Z75D3
\paragraph{Z75D3. Counting neighbors with the class-sum weights.}
Take $M>0$, a prime $P$ as in Z75D2 for $m=M$, and put $Q=\bar P$.
Let $L$ range over the finitely generated self-dual Hermitian
$O$-lattices of the source's $V(M)$ in $E^n$. They agree with
$O^n$ at level $M$, meaning that some positive integer prime to $M$
carries either lattice into the other by multiplication. Denote by
$G=D_H(M)$ the source's rational unitary group with level condition.
A $P$-neighbor $L'$ lies in $V(M)$ and satisfies
$PL\subset L'$ and $[L:L\cap L']=q$.

Such neighbors are in bijection with the hyperplanes in $L/PL$.
Indeed self-duality and $PL\subset L'$ also give $QL'\subset L$,
by moving scalars across the Hermitian pairing. At a finite prime
different from $P,Q$, the local lattices therefore coincide.
At $P$, $L'_P\subset L_P$ contains $P L_P$, and the index condition
makes its reduction a hyperplane. At $Q$, $L'_Q$ contains $L_Q$.
Under the local Hermitian duality between the $P$- and $Q$-spaces,
$L'_Q$ must be the dual of $L'_P$. Conversely take the inverse
image of any such hyperplane at $P$, its dual at $Q$, and the
original lattice elsewhere. These prescriptions define a unique
lattice. For instance, all the prescribed lattices lie between
$qL$ and $q^{-1}L$ locally, and the primary decomposition of the
finite module $q^{-1}L/qL$ imposes the prescriptions independently.
This lattice is self-dual, has the required index, and agrees
with $L$ at level $M$ using the multiplier $q$. Thus it belongs
to $V(M)$. Locally at the discrete valuation ring $O_P$ the passage
to a hyperplane multiplies one basis vector by a uniformizer,
so the dual enlargement at $Q$ also has index $q$. Consequently
$L'$ is a $P$-neighbor of $L$ exactly when $L$ is a $Q$-neighbor
of $L'$. These descriptions commute with $G$.

The rank of $L$ is $n$ by agreement with $O^n$, and its fiber over
$P$ has dimension $n$ by freeness over $O_P$. There are
$(q^n-1)/(q-1)$ hyperplanes, one for each line of nonzero linear
forms up to scalar. More generally if $u:X\to L$ is $O$-linear
and its reduction has rank $s$, precisely $(q^{n-s}-1)/(q-1)$
neighbors contain $u(X)$, since their hyperplanes must contain its
reduction.

For a $G$-class function $f$ use the finite definite class set
and finite stabilizers of the native class sum. For representatives
$L_i,L_j$, let $m_P(i,j)$ be the number of $P$-neighbors of $L_i$
in the class of $L_j$, and $s_i=|\operatorname{Stab}_G(L_i)|$.
The elements $g\in G$ with $gL_j$ a $P$-neighbor of $L_i$ are
counted $s_j$ for each of these neighbors. Inversion counts the
same set by the reversed $Q$-neighbors, giving
$m_P(i,j)s_j=m_Q(j,i)s_i$. Interchanging the two finite sums gives
\[
 (q-1)\sum_i\frac1{s_i}\sum_{L'\text{ a }P\text{-neighbor of }L_i}f(L')
   =(q^n-1)\sum_j\frac{f(L_j)}{s_j}.
\]
The weight on the left belongs to the class of the starting lattice.
% END Z75D3

% BEGIN Z75D4
\paragraph{Z75D4. The two lattice families and their residue counts.}
Now let $P$ satisfy Z75D2 for $m=RM$, with $M,R>0$, and $d>0$.
Let $X\subset E^d$ be a finitely generated $O$-lattice agreeing
with $O^d$ at level $RM$. Put $Q=\bar P$. Both $X/PX$ and $X/QX$
have dimension $d$ over their residue fields of size $q$, by the
rank and local freeness argument of Z75D3. Set $A=(q^d-1)/(q-1)$.
Choose lifts $e_i\in X$ of nonzero representatives of every line of
$X/PX$, one per line. Independently enumerate by the same index set
the $A$ inverse images $B_i\subset X$ of hyperplanes in $X/QX$.

Choose $t\in P^{-1}\setminus O$. Such elements exist because $P$ is a
proper invertible ideal. The valuation of $t$ is $-1$ at $P$ and is
nonnegative elsewhere. Set
\[
 X_{A,i}=X+Ote_i,\qquad X_{B,i}=B_i.
\]
The quotient $X_{A,i}/X$ is generated by $te_i$, is killed by $P$,
and is nonzero since $e_i$ is nonzero modulo $PX$ and $t$ has a
simple pole at $P$. Thus it has size $q$. Likewise $X/B_i$ has
size $q$. These lattices agree with $X$ at level $RM$ using $q$.
Write $I(X)=[X:X\cap O^d]/[O^d:X\cap O^d]$ for the source's
positive generalized index. Multiplicativity of finite indices gives
\[
 I(X_{A,i})=qI(X),\qquad I(X_{B,i})=I(X)/q.
\]

For an $O$-linear map $u:X\to L$ as in Z75D3, let $s$ be its
rank on reduction modulo $P$. The number $a$ of indices for which
$t u(e_i)\in L$ counts exactly the lines of its kernel. The number
$b$ of $P$-neighbors of $L$ containing $u(X)$ counts hyperplanes
through its image. In particular
\[
 ((q-1)a+1)q^n=q^{d-s}q^n
               =q^{n-s}q^d=((q-1)b+1)q^d .
\]
This is the finite-field equality used on the integral branch of
the pointwise Hecke calculation. For the other family the native
inner sum uses $q^d$ lifts $\xi_j$ of $X/PX$. Any additive character
of $X$ trivial on $PX$ sums to $q^d$ on these lifts if it is trivial
on $X$, and to zero otherwise, by translating the sum by an element
whose character value differs from one.
Thus in degree one there is one index of each lattice type and $q$
inner terms. In degree two there are $q+1$ indices of each type
and $q^2$ inner terms per index.
% END Z75D4

% BEGIN Z75D5
\paragraph{Z75D5. The exact contraction at a positive degree.}
Fix an admitted family $\phi$ at levels $M,R$, even $n=2k$, and a
positive degree $d$ with $4d<n$. Use the prime $q,P$ of Z75D2 at
level $RM$. Write $D_X(Z)$ for the class-sum difference of Z75D1
applied to the native twist of $\phi$ by $X$, for $X$ as in Z75D4.
Let $J(Z)=\prod_{\tau\in\Psi}\det(\operatorname{Im}Z_\tau)>0$.
Suppose that a constant $C$ bounds
\[
 F(X,Z)=I(X)^k J(Z)^k |D_X(Z)|
\]
throughout this entire family of lattices and periods. This uniform
bound is the use of the lower-degree Siegel--Weil identity in the
native successor proof. Also suppose the weighted Hecke identity
holds on this same family, with the lattices of Z75D4 and Hermitian
translation matrices $\beta_{ij}$. In the notation of P75H3 this reads
\[
\begin{split}
 (q^d-1)(q^{2k}+q^d)D_X(Z)
   =(q-1)\left(q^{2k}\sum_{i=1}^{A}D_{X_{A,i}}(Z)
   +\sum_{i=1}^{A}\sum_{j=1}^{q^d}
           D_{X_{B,i}}(Z+\beta_{ij})\right).
\end{split}
\]
The translations are taken at their images under all embeddings.
Then $D_\phi(Z)=0$ for every $Z$ of degree $d$.

Indeed let $S=\sup_{X,Z}F(X,Z)$. It is finite by the uniform bound
and is nonnegative. Hermitian translations preserve $J$. After
multiplying the Hecke identity by $I(X)^kJ(Z)^k$ and taking
absolute values, the contribution of each term of the first sum
inside parentheses is at most $q^{2k}q^{-k}S=q^k S$, by the first
index formula of Z75D4. Each term of the second sum is at most
$q^k S$ by the other index formula. There are $A$ and $Aq^d$
terms, respectively. Using $(q-1)A=q^d-1$ yields
\[
 S\le\frac{q^k(1+q^d)}{q^{2k}+q^d}\,S.
\]
The denominator exceeds the numerator by
$(q^k-1)(q^k-q^d)>0$. Hence $S=0$. In particular $D_{O^d}=0$.

For completeness the standard lattice twist gives $\phi_\Lambda(x)$
at every $\Lambda\in V(M)$. Its support condition tests whether all
columns of $Rx$ lie in $\Lambda$. If this fails, admission gives
$\phi_\Lambda(x)=0$. Otherwise the twist evaluates at $sx$, with
$s\equiv1\pmod{R^2}$ and $sO^d\subset O^d$ (one may take $s=1$).
For every allowed choice, $(s-1)x$ is $R$ times a matrix with
columns in $\Lambda$, so periodicity gives $\phi_\Lambda(sx)=
\phi_\Lambda(x)$. This proves the passage to $D_\phi$.
The required uniformity in $X$ and $Z$ in the first hypothesis
continues to require proof for a successor degree.
% END Z75D5

% BEGIN Z75D6
\paragraph{Z75D6. A single level for all rationally moved twists.}
For fixed $E$ and degree $d>0$, there is an integer $c>0$
such that every $X$ of Z75D4 admits a rational invertible matrix
$a\in\mathrm{GL}_d(E)$ with
\[
 cO^d\subset aX\subset O^d,\qquad
 I(X)\le |N_{E/\mathbb Q}\det a|.
\]
Moreover the family $\psi_\Lambda(x)=(\mathrm{twist}_R(X)\phi)_\Lambda(xa)$
is admitted at $(M,Rc)$, simultaneously for all $X$ with this same $c$.

The ideal class group of $O$ is finite. Choose a nonzero integral
ideal in each class and choose $c>0$ lying in every one of these
finitely many ideals. Every full lattice $X$ is isomorphic as an
$O$-module to a direct sum of $d$ ideals from this list. Indeed a
nonzero coordinate projection has image a fractional ideal, which
is projective over the Dedekind domain $O$. Splitting off the image
and iterating on the kernel gives a sum of fractional ideals, each
isomorphic by scaling to its chosen representative. Extending the
module isomorphism to $E$ gives the asserted $a$. The generalized
index is the covolume of $O^d$ divided by the covolume of $X$ in
Minkowski space. Hence $I(aX)=I(X)/|N_{E/\mathbb Q}\det a|$,
and $aX\subset O^d$ gives $I(aX)\le1$.

Take $s\equiv1\pmod{R^2}$ with $sO^d\subset X$, obtained by choosing
a multiple of a level-agreement multiplier prime to $R$.
At a fixed $\Lambda$, the twist at $xa$ vanishes unless
$RxaX\subset\Lambda$, and on this set equals $\phi_\Lambda(sxa)$.
The value is independent of the eligible $s$: if $s'$ is also
eligible, $Rsxa$ and $Rs'xa$ have columns in $\Lambda$. Periodicity
gives the same value at $sxa$ and $ss'xa$ since $s'-1$ is a
multiple of $R^2$, and likewise at $s'xa$ and $ss'xa$.
On the support, $cO^d\subset aX$ forces the columns of $Rcx$ into
$\Lambda$. Translating $x$ by $Rcy$ with $yO^d\subset\Lambda$
preserves the support test since $aX\subset O^d$, and preserves the
value since $saO^d\subset aX\subset O^d$. Left multiplication
by a unitary element carrying $\Lambda$ to $\Lambda'$ transports
the support test and preserves the value by equivariance of $\phi$.

For the locality clause let $\Lambda,\Lambda'$ agree via a multiplier
$u$ coprime to $RcM$, and let $Rcx$ have its columns in both.
If $w=Rxa\xi\in\Lambda$ for $\xi\in X$, then $uw\in\Lambda'$.
Also $cw\in\Lambda'$ since $a\xi\in O^d$. B\'ezout gives
$w\in\Lambda'$, and the converse follows by symmetry. The two support
tests therefore agree. On that support $Rsxa$ has columns in both
lattices, so the original locality clause at $RM$ gives equal test
values. This proves every clause of admission at the fixed level.

Consequently the uniform bound in Z75D5 follows from the fixed-level
estimate $J(W)^k|D_\psi(W)|\le C\sup_{\Lambda,x}|\psi_\Lambda(x)|$
for every $(M,Rc)$-admitted family, under the native parabolic Weil
transformation. Indeed take $Z'_\tau=a_\tau Z_\tau a_\tau^*$.
For the rational unitary block matrix $\operatorname{diag}
(a,(a^*)^{-1})$, its zero-variable automorphy scalar
$\prod_\tau\det((a_\tau^*)^{-1})^n$ occurs both in the Weil action on
the test and in the period transformation, and cancels. Thus
$D_\psi(Z')=D_X(Z)$. Also
$J(Z')=|N_{E/\mathbb Q}\det a|J(Z)$ by the CM-type product.
The supremum of $|\psi|$ is bounded by that of the original $|\phi|$,
which is finite: on each of finitely many lattice classes the
admission conditions restrict the test to finitely many residues,
and equivariance transports the bound. The claimed bound for
$I(X)^kJ(Z)^k|D_X(Z)|$ follows, uniformly in $X$ and $Z$.
The fixed-level estimate still requires the lower-degree identity
and the analytic argument at the cusps.
% END Z75D6
